% ===================================================================== Chapter 8
\chapter{Baseline CFD Results}\label{ch:cfdres}
All values in this chapter belong to the official baseline P00 on the medium mesh (159,840~cells) and are summarised in
Table~\ref{tab:cfdres}. They are simulation results of the baseline problem.

%% generated by gen_tables.py - RE-ANALYSIS 2026
\begin{table}[htbp]
\centering
\small
\caption{Baseline CFD results (medium mesh, converged, official baseline P00)}
\label{tab:cfdres}
\begin{tabular}{>{\raggedright\arraybackslash}p{0.50\linewidth}rl}
\toprule
Quantity & Value & Unit \\
\midrule
Cells (fluid + solid) & 159,840 & -- \\
Reynolds number, inlet / outlet & 29,958 / 25,545 & -- \\
Pressure drop (area-weighted static, inlet $-$ outlet) & 438.13 & Pa \\
Outlet bulk temperature (mass-weighted) & 368.93 & K \\
Heat-transfer rate through the heated wall & 602.755 & W \\
Maximum solid temperature (outer-wall facet) & 562.58 & K \\
Maximum solid temperature (hottest cell centre) & 562.13 & K \\
Through-wall $\Delta T$ at mid-span ($z$ = 300 mm) & 7.581 & K \\
Fully developed Nusselt number ($x/D$ 18--29) & 54.17 & -- \\
Fully developed Darcy friction factor ($x/D$ 18--29) & 0.02163 & -- \\
$y^+$ minimum / area mean / maximum & 0.208 / 0.241 / 0.585 & -- \\
Mass imbalance / energy imbalance & $7.0\times10^{-13}$ / $3.1\times10^{-11}$ & \% \\
\bottomrule
\end{tabular}
\par\smallskip\parbox{\linewidth}{\footnotesize\raggedright Source: \texttt{MASTER\_PROJECT\_DATA.csv} rows M033--M048 (VERIFIED from the Fluent flux and surface-integral reports, except M035 REPORTED and M041--M043 CROSS-CHECKED). $Q$ = 602.755~W is 8000~W/m$^2$ times the inscribed 48-gon area; the same flux on the true cylinder is 603.186~W.}
\end{table}


\section{Convergence}
The solution converged at iteration 600 and was confirmed at iteration 700 (Figure~\ref{fig:conv}). The evaluation at iteration 500 was
refused because the 200-iteration window still contained the transient caused by the switch to second-order schemes (outlet temperature
drift 0.42~K). The residuals fell to round-off level (continuity $9.0\times10^{-11}$, energy $1.6\times10^{-14}$); the evidence of convergence,
however, is the physical monitors, which are flat to within $10^{-5}$ K. The solid temperature rose monotonically during the energy stage,
without overshoot.

\begin{figure}[htbp]
\centering
\includegraphics[width=0.90\linewidth]{fig16_convergence_nt.png}
\caption[Convergence history of the baseline run]{Convergence history of the baseline run: (a) scaled residuals, (b) monitored temperatures, (c) pressure drop and interface heat rate.
Stage boundaries at iterations 150 (energy on) and 300 (second-order schemes).}
\label{fig:conv}
\src{\provcfd}
\end{figure}

\section{Mass Conservation}
Inlet and outlet mass flow are both 8.662155~g/s; the Fluent flux report gives a mass imbalance of $7.0\times10^{-13}$\,\% of the flow. The mass
flow is 0.28\,\% below the analytical 8.686~g/s, which is entirely explained by the inscribed-polygon inlet area (area ratio 0.997147) and the
gas constant used by Fluent. No reverse-flow cell or outlet face occurred.

\section{Energy Conservation}
The heat entering the outer wall (602.755~W) crosses the interface to within $9\times10^{-10}$ W and leaves with the air: the enthalpy flux rises
from 16.107 to 618.862~W. The net energy imbalance over all boundaries is $3.1\times10^{-11}$\,\% of $Q$. The 0.071\,\% difference from the
analytical 603.186~W is exactly the ratio of the 48-gon lateral area to the cylinder area.

\section{Velocity Field}
The uniform inlet profile develops into a turbulent profile (Figure~\ref{fig:vel}). The ratio of maximum to bulk velocity reaches 98\,\% of its
exit value at $x/D \approx 17.8$. Because heating lowers the density by about 19\,\% along the duct, the flow accelerates: the area-mean
velocity rises from 23.500~m/s to 28.894~m/s and the maximum velocity, 34.99~m/s, occurs on the outlet centreline. Radial and swirl
velocities stay below 0.54~m/s and 0.013~m/s; secondary flow is negligible.

\begin{figure}[htbp]
\centering
\begin{subfigure}[t]{0.49\linewidth}\centering\includegraphics[width=\linewidth,height=52mm,keepaspectratio]{fig01_velocity_contour_nt.png}\caption{Axial velocity (radial scale exaggerated) [m/s]}\end{subfigure}\hfill
\begin{subfigure}[t]{0.49\linewidth}\centering\includegraphics[width=\linewidth,height=52mm,keepaspectratio]{fig13_velocity_profiles_nt.png}\caption{Normalised axial-velocity profiles $w/w_b$}\end{subfigure}
\caption[Velocity field of the baseline]{Velocity field of the baseline: development from the uniform inlet profile to a turbulent profile, and acceleration of the heated air.}
\label{fig:vel}
\src{\provcfd{} The 1/7 power law in (b) is a textbook reference shape only.}
\end{figure}

\section{Pressure Field}
The static pressure is almost uniform over each cross-section and falls monotonically from 438.13~Pa (area average at the inlet) to 0~Pa
gauge at the outlet (Figure~\ref{fig:press}). The gradient is steepest in the entry region, where the boundary layer is thin, and steepens
again downstream as heating accelerates the air. An axial momentum balance on the exported field decomposes the drop into wall shear
(242.87~Pa), one-dimensional acceleration (149.31~Pa) and momentum-flux profile development (46.15~Pa), closing to $-$0.045\,\%.

\begin{figure}[htbp]
\centering
\begin{subfigure}[t]{0.52\linewidth}\centering\includegraphics[width=\linewidth,height=52mm,keepaspectratio]{fig04_pressure_contour_nt.png}\caption{Static gauge pressure [Pa]}\end{subfigure}\hfill
\begin{subfigure}[t]{0.46\linewidth}\centering\includegraphics[width=\linewidth,height=52mm,keepaspectratio]{fig10_pressure_axial_nt.png}\caption{Area-averaged static pressure along the duct}\end{subfigure}
\caption[Pressure field of the baseline]{Pressure field of the baseline: inlet-to-outlet static pressure drop of 438.13~Pa.}
\label{fig:press}
\src{\provcfd}
\end{figure}

\section{Temperature Field}
The air heats almost linearly along the duct, from 300~K to a mass-weighted outlet temperature of 368.93~K (Figures~\ref{fig:temp}
and~\ref{fig:taxial}a). The heated layer is thin: at the outlet a hot wall layer surrounds a core that is still close to the bulk temperature.
The solid ranges from 423.97~K (inlet end) to 562.58~K (medium mesh) at the outer-wall facet near the adiabatic outlet end; the hottest cell centre is at
562.13~K and the volume-mean solid temperature is 525.48~K. The through-wall temperature difference is 7.581~K at mid-span and 7.361~K at
$z$ = 570~mm, but rises to 15.1~K in the first axial slab, where the cold air and thin thermal boundary layer draw heat axially through the
wall towards the inlet (Figure~\ref{fig:taxial}b).

\begin{figure}[htbp]
\centering
\begin{subfigure}[t]{0.40\linewidth}\centering\includegraphics[width=\linewidth,height=52mm,keepaspectratio]{fig05_fluid_temperature_contour_nt.png}\caption{Air temperature [K]}\end{subfigure}\hfill
\begin{subfigure}[t]{0.58\linewidth}\centering\includegraphics[width=\linewidth,height=52mm,keepaspectratio]{fig06_solid_temperature_contour_nt.png}\caption{Inconel 718 wall temperature [K]}\end{subfigure}
\caption{Temperature field of the baseline in the meridional plane (radial scale exaggerated) and in three wall cross-sections.}
\label{fig:temp}
\src{\provcfd}
\end{figure}

\begin{figure}[htbp]
\centering
\begin{subfigure}[t]{0.48\linewidth}\centering\includegraphics[width=\linewidth,height=52mm,keepaspectratio]{fig11_bulk_temperature_axial_nt.png}\caption{Bulk and wall temperatures, CFD and analytical}\end{subfigure}\hfill
\begin{subfigure}[t]{0.50\linewidth}\centering\includegraphics[width=\linewidth,height=52mm,keepaspectratio]{fig12_through_wall_temperature_nt.png}\caption{Radial profile and through-wall $\Delta T$ along the duct}\end{subfigure}
\caption[Axial temperature development and through-wall temperature difference]{Axial temperature development and through-wall temperature difference. Dashed and dotted lines are the Section 2 analytical estimates.}
\label{fig:taxial}
\src{\provcfd}
\end{figure}

\section{Heat Transfer}
Heat enters uniformly at 8000~W/m$^2$ on the outside but does not leave uniformly at the bore (Figure~\ref{fig:flux}a). In the first slab
the bore flux reaches 38,471~W/m$^2$, 2.4 times the one-dimensional 16,000~W/m$^2$, because the heat-transfer coefficient is very high in the
thermal entrance region; it recovers to within 1\,\% of 16,000~W/m$^2$ by $z \approx$ 250~mm. Figure~\ref{fig:flux}b shows that almost the
whole wall-to-bulk temperature difference lies in the thin air layer next to the wall: the 10~mm of Inconel carries only a few kelvin.
This is the Section 2 finding that the film carries about 97\,\% of the thermal resistance, seen directly in the solution.

\begin{figure}[htbp]
\centering
\begin{subfigure}[t]{0.54\linewidth}\centering\includegraphics[width=\linewidth,height=52mm,keepaspectratio]{fig08_wall_heat_flux_nt.png}\caption{Wall heat flux [W/m$^2$]}\end{subfigure}\hfill
\begin{subfigure}[t]{0.44\linewidth}\centering\includegraphics[width=\linewidth,height=52mm,keepaspectratio]{fig14_radial_temperature_nt.png}\caption{Radial temperature in air and wall}\end{subfigure}
\caption{Heat-flux distribution at the interface and radial temperature distribution at three stations.}
\label{fig:flux}
\src{\provcfd}
\end{figure}

\section{Nusselt Number}
Over the fully developed window $x/D$ = 18--29 the mean Nusselt number is 54.17 (Figure~\ref{fig:nuf}a); the corresponding film coefficient
near the outlet is 83.6~W/(m$^2$\,K). The value lies between the property-corrected correlation (51.0 at the CFD conditions) and the
constant-property correlations (Gnielinski 63.5, Dittus--Boelter 68.6), and inside the expected band of 44--67. The CFD implies a
property-ratio exponent of about 0.36 instead of the 0.5 assumed in the analytical model.

\begin{figure}[htbp]
\centering
\includegraphics[width=0.78\linewidth]{fig15_nusselt_friction_nt.png}
\caption[Local Nusselt number (a) and Darcy friction factor (b) along the duct compared with correlations]{Local Nusselt number (a) and Darcy friction factor (b) along the duct compared with correlations; the shaded band is the fully developed window $x/D$ = 18--29.}
\label{fig:nuf}
\src{\provcfd}
\end{figure}

\section{Friction Factor}
The fully developed Darcy friction factor, $f = 8\tau_w\rho_b/G^2$ over the same window, is 0.02163, 10.4\,\% below the Petukhov value of
0.0241 (Figure~\ref{fig:nuf}b). This marginally misses the $\pm$10\,\% requirement set in Section 1 (PR-04), which is recorded rather than
re-interpreted. The mesh study (Chapter~\ref{ch:mi}) traces the difference mainly to gas heating: the variable-property effect reduces the
friction factor by a factor of 0.910 beyond the constant-property Reynolds-number change, which the constant-property correlation omits.

\section{y+}\label{sec:yplus}
On the conjugate wall $y^+$ ranges from 0.208 to 0.585 with an area mean of 0.241 (Figure~\ref{fig:yplus}); 100\,\% of the area has
$y^+ \le 1$ and 99.07\,\% has $y^+ \le 0.5$. The maximum occurs in the first slab, where the wall shear is highest. Wall shear and wall heat
flux are therefore resolved, not modelled by a wall function.

\begin{figure}[htbp]
\centering
\includegraphics[width=0.85\linewidth]{fig09_yplus_nt.png}
\caption[Wall $y^+$ on the fluid--solid interface]{Wall $y^+$ on the fluid--solid interface: (a) unwrapped distribution, (b) along the duct, (c) area-weighted histogram.}
\label{fig:yplus}
\src{\provcfd{} The four-fold circumferential pattern in (a) is geometric (thinner first cells at the O-grid diagonals).}
\end{figure}

\section{Analytical vs CFD Comparison}
Table~\ref{tab:anvscfd} compares the Section 2 analytical estimates with the CFD. The two are different modelling levels and neither was
adjusted towards the other. The conservation-based quantities agree to within the exact faceting effect. The maximum solid temperature is
19.1~K lower in the CFD (562.58~K against the analytical 581.71~K): 15.0~K of the difference comes from the higher CFD film coefficient and 4.2~K from axial
conduction towards the adiabatic outlet end, which the one-dimensional model cannot represent. Mesh refinement moves the CFD further from
the analytical value (Chapter~\ref{ch:mi}), so the gap is a model difference, not a mesh artefact. Table~\ref{tab:appB} lists the analytical
pressure-drop components.

%% generated by gen_tables.py - RE-ANALYSIS 2026
\begin{table}[htbp]
\centering
\footnotesize
\caption{Section 2 analytical estimates compared with the baseline CFD (different modelling levels)}
\label{tab:anvscfd}
\begin{tabular}{>{\raggedright\arraybackslash}p{0.25\linewidth}rrr>{\raggedright\arraybackslash}p{0.30\linewidth}}
\toprule
Quantity & Analytical & CFD medium & Difference & Explanation \\
\midrule
$Q$ [W] & 603.186 & 602.755 & $-$0.071\,\% & exactly the 48-gon lateral-area ratio \\
$T_{out}$ [K] & 368.85 & 368.93 & +0.08 K & faceting of the CFD section \\
$\Delta p$ [Pa] & 441.53 & 438.13 & $-$0.77\,\% & friction 7.7\,\% lower, profile-development term higher; partly compensating \\
$T_{max}$ solid [K] & 581.71 & 562.58 & $-$19.1 K & higher CFD film coefficient ($-$15.0 K) and axial conduction ($-$4.2 K) \\
$\Delta T_{wall}$ [K] & 7.285 (exit) & 7.361 ($z$ = 570 mm) & +1.0\,\% & 1-D conduction confirmed away from the ends \\
$Nu_{fd}$ & 49.6 & 54.17 & +9.3\,\% & CFD implies property exponent $n \approx 0.36$ instead of 0.5 \\
$f_{fd}$ & 0.0241 & 0.02163 & $-$10.4\,\% & gas-heating (variable-property) effect, F-029 \\
\bottomrule
\end{tabular}
\par\smallskip\parbox{\linewidth}{\footnotesize\raggedright Sources: \texttt{MASTER\_PROJECT\_DATA.csv} rows M036--M043 and M071--M075; explanations from \texttt{CFD\_VS\_ANALYTICAL.md} and \texttt{PROJECT\_STATE.md} \S5b. Neither level was adjusted to match the other.}
\end{table}


